\paragraph{A separating source must persist.}
Canonical realizability imposes a condition across different Boolean cubes.
Write $f(A)$ for the value at the characteristic assignment of $A$.

\begin{proposition}[A stable separating source]
\label{core:prop:canonical-separating-source}
If $D_f$ is realized by a bichain term, including observer union, then
for every $f(A)=1$ and $f(B)=0$ there is one $j\in A\setminus B$
essential on every cube $[L,U]$ with $L\subseteq B$, $A\subseteq U$,
and $L\subseteq U$.
\end{proposition}
\begin{proof}
Follow a separating path in a realizing term. At a join choose a child
true at $A$; both children are false at $B$. At a meet choose a child
false at $B$; both are true at $A$. An observer union giving endpoints
at both assignments has separating children, so choose either one.
The resulting leaf is a variable $j\in A\setminus B$, chosen before
the cube.

Mark $j$ on any displayed cube and leave its other free coordinates
neutral. The relation $x=c\Rightarrow y\ne b$ is preserved
coordinatewise by all three bichain operations. Along the chosen path, a meet sibling cannot be the false
endpoint: it is true at $A$, whose true coordinates are not fixed false
on the cube. Dually a join sibling cannot be the true endpoint, since
it is false at $B$ and every coordinate fixed true lies in $B$.
The marked value therefore survives to the root, also through observer
union. Canonical semantics makes $j$ essential on that cube.
\end{proof}

This condition excludes a five-variable operation even though fixing
any one Boolean coordinate gives an exact read-once realizer.

\begin{theorem}[A Boolean-cofactor obstruction]
\label{core:thm:canonical-cofactor-obstruction}
Put
\[
 F(a,b,c,d,z)=ab\lor cd\lor zad\lor zbc.
\]
Its canonical extension $D_F$ is not a bichain term operation, even
with observer union. Each of its ten one-variable Boolean cofactors
has a pure read-once term realizing its canonical extension.
\end{theorem}
\begin{proof}
Take $A=\{a,d,z\}$ and $B=\{b,c\}$. Then $F(A)=1$ and $F(B)=0$.
On the compatible cube with $c=1,b=0$, the function is the projection
$d$; on the compatible cube with $b=1,c=0$, it is the projection $a$.
Their essential sets are disjoint, contradicting
\cref{core:prop:canonical-separating-source}.

The cofactor identities below have read-once right sides:
\[
\begin{array}{c|c}
 a=0&c\wedge(d\vee(z\wedge b))\\
 a=1&b\vee(d\wedge(c\vee z))\\
 z=0&(a\wedge b)\vee(c\wedge d)\\
 z=1&(a\vee c)\wedge(b\vee d).
\end{array}
\]
The permutations $(a\ b)(c\ d)$ and $(a\ c)(b\ d)$ preserve $F$
and act transitively on its first four coordinates, giving the other
six cofactors. A pure read-once formula has exact canonical semantics.
Stable outputs are determined by the cube corners. At a possible output,
choose a child carrying a marked source; its sibling is not a stable
absorber. The disjoint supports let a pivotal edge in that child combine
with a Boolean endpoint assignment for its sibling. Conversely, an
essential marked coordinate forces possible in every Boolean realizer.
Induction therefore gives equality with $D_f$. Apply this to the
displayed cofactors and their permutations.
\end{proof}

The same operation also has exact realizers on every two-source marked
fragment, even with arbitrary neutral inputs elsewhere.

\begin{corollary}[Pairwise marked-source locality]
\label{core:cor:canonical-pair-locality}
For each set $J$ of at most two coordinates of $F$, there is one pure
term $t_J$ such that
\[
 t_J(v)=D_F(v)\qquad
 \text{whenever }\{i:v_i=\pv\}\subseteq J.
\]
\end{corollary}
\begin{proof}
Let $Z_J$ be the semantic budget of $F$ restricted to marks in $J$.
Its direct graph is \eqref{core:eq:direct-marked-edge}. We show that
$F$ respects every direct edge. Its minterms are $ab,cd,zad,zbc$, and
its maximal false facets are $bc,ad,zac,zbd$. Every minterm-minus-facet
has size one, except $zad\setminus bc$ and $zbc\setminus ad$, which
have size three.

If a direct edge crossed from true to false, shrink its source to a
minterm $A$ and enlarge its target to a false facet $B$. The same row
witnesses remain valid, so $A\setminus B\subseteq J$ and this difference
must be a singleton $\{j\}$. But $j$ is essential on every compatible
cube $L\subseteq B$, $U\supseteq A$: the assignment
$L\cup(A\setminus\{j\})\subseteq B$ is false, and adding $j$ makes
it true. No row of $Z_J$ can witness the required inessentiality.

By \cref{core:thm:two-marked-budgets}, $F$ therefore has a pure formula
safe on all of $Z_J$. Its endpoint constants simplify away because $F$
is nonconstant, giving the required $t_J$.
Stable outputs agree with the cube endpoints, and an essential marked
coordinate forces $\pv$ in every Boolean realizer. For the remaining
inputs, a term returns $\pv$ exactly when it does so with just one of
the marked coordinates retained: replacing other $\pv$ inputs by
$\nv$ leaves all stable-endpoint tests unchanged, and each gate
passes the disjunction of its unabsorbed marked-source flags.
Every such single-mark input is in $Z_J$, where the formula is safe.
Thus all inputs in the stated fragment agree with $D_F$.
\end{proof}
